\documentclass[10pt]{beamer}
\usetheme{default}
\setbeamercovered{invisible}
\setbeamertemplate{navigation symbols}{}
\setbeamertemplate{footline}{
    \flushright{\hfill \insertframenumber{}/\inserttotalframenumber}
}

\usepackage{listings}

% User-defined colors.
\definecolor{DarkGreen}{rgb}{0, .5, 0}
\definecolor{DarkBlue}{rgb}{0, 0, .5}
\definecolor{DarkRed}{rgb}{.5, 0, 0}
\definecolor{LightGray}{rgb}{.95, .95, .95}
\definecolor{White}{rgb}{1.0,1.0,1.0}
\definecolor{darkblue}{rgb}{0,0,0.9}
\definecolor{darkred}{rgb}{0.8,0,0}
\definecolor{darkgreen}{rgb}{0.0,0.85,0}

% Settings for listing class.
\lstset{
  language=C++,                        % The default language
  basicstyle=\small\ttfamily,          % The basic style
  backgroundcolor=\color{White},       % Set listing background
  keywordstyle=\color{DarkBlue}\bfseries, % Set keyword style
  commentstyle=\color{DarkGreen}\itshape, % Set comment style
  stringstyle=\color{DarkRed}, % Set string constant style
  extendedchars=true % Allow extended characters
  breaklines=true,
  basewidth={0.5em,0.4em},
  fontadjust=true,
  linewidth=\textwidth,
  breakatwhitespace=true,
  showstringspaces=false,
  lineskip=0ex, %  frame=single
}

\begin{document}
    \title{Plugins \protect\\ Generic programming}
    \author{Pasquale Claudio Africa}
    \date{05/06/2020}

\begin{frame}[plain, noframenumbering]
    \maketitle
\end{frame}

\begin{frame}{Exercise 1: quadrature rules with plugins}
\begin{itemize}
\item Implement a code that enables the user to compute the integral of a scalar function by selecting the quadrature rule as the name of a dynamically loadable object;
\item The dynamically loadable object should define a function with the following signature:
\texttt{double integrate(std::function<double (double)>, double a, double b)}.
\item Implement plugins for midpoint, trapezoidal and Simpson's rule.
\item Implement a plugin for quadrature with adaptive refinement.
\end{itemize}
\end{frame}


\begin{frame}{Exercise 2: 1D interpolator using generic programming}
This exercise (inspired by \texttt{Examples/src/RKFSolver}) shows an example of a generic piecewise linear 1D interpolator, with an application to
vectors of couples of values (representing the interpolation nodes and the corresponding values) and to two separate vectors (one for the nodes and one for the values).

It is rather generic and accepts as an input iterators to any container with bi-directional iterators. Bi-directionality is indeed only needed for extrapolation on the right.
\end{frame}

\begin{frame}{Thank you all and have fun with \texttt{C++}!}
\centering
\includegraphics[width=0.75\textwidth]{meme_c++.png}\\
\url{https://www.reddit.com/r/cpp/comments/4m4pep/fun_c_syntax/}
\end{frame}

\end{document}
